\documentclass[12pt, a4paper]{article}

% ── PACKAGES ──────────────────────────────────────────────────
\usepackage[utf8]{inputenc}
\usepackage[T1]{fontenc}
\usepackage{geometry}
\geometry{margin=2.5cm}
\usepackage{hyperref}
\hypersetup{
    colorlinks=true,
    linkcolor=blue,
    urlcolor=blue,
    citecolor=blue,
    pdftitle={Tazemetostat as Maintenance Therapy
              After Neoadjuvant Chemotherapy in
              EZH2-High TNBC},
    pdfauthor={Eric Robert Lawson},
    pdfkeywords={tazemetostat, maintenance therapy,
                 TNBC, EZH2, EZH2 paradox, late relapse,
                 neoadjuvant chemotherapy, breast cancer,
                 OrganismCore, attractor geometry,
                 GSE25066, post-chemotherapy window}
}
\usepackage{amsmath}
\usepackage{booktabs}
\usepackage{longtable}
\usepackage{array}
\usepackage{parskip}
\usepackage{microtype}
\usepackage{authblk}
\usepackage{titlesec}
\usepackage{fancyhdr}
\usepackage{xcolor}
\usepackage{float}
\usepackage{enumitem}
\usepackage{setspace}

% ── COLOURS ───────────────────────────────────────────────────
\definecolor{repobox}{RGB}{240,247,255}
\definecolor{urgentbox}{RGB}{255,235,235}
\definecolor{lockbox}{RGB}{245,255,245}
\definecolor{mechanismbox}{RGB}{250,245,255}
\definecolor{paradoxbox}{RGB}{255,250,235}

% ── HEADER / FOOTER ───────────────────────────────────────────
\pagestyle{fancy}
\fancyhf{}
\rhead{\small OrganismCore \textbar{} 2026-03-06}
\lhead{\small CS-LIT-17: Tazemetostat Maintenance
       in EZH2-High TNBC}
\rfoot{\small Page \thepage}
\renewcommand{\headrulewidth}{0.4pt}

% ── SECTION FORMATTING ────────────────────────────────────────
\titleformat{\section}
  {\large\bfseries}{\thesection}{1em}{}[\titlerule]
\titleformat{\subsection}
  {\normalsize\bfseries}{\thesubsection}{1em}{}
\titleformat{\subsubsection}
  {\small\bfseries}{\thesubsubsection}{1em}{}

% ══════════════════════════════════════════════════════════════
% TITLE
% ══════════════════════════════════════════════════════════════
\title{
    \vspace{-1em}
    \textbf{Tazemetostat as Maintenance Therapy
    After Neoadjuvant Chemotherapy in EZH2-High
    Triple-Negative Breast Cancer:}\\[0.5em]
    \large Exploiting the Post-Chemotherapy Window
    to Close the Late-Relapse Gap\\[0.4em]
    \normalsize \textit{OrganismCore Framework ---
    Document CS-LIT-17}\\[0.3em]
    \small Derived from Waddington Landscape Attractor
    Geometry and the EZH2 Paradox\\[0.3em]
    \small Companion to CS-LIT-16:
    Tazemetostat $\rightarrow$ Fulvestrant
    in EZH2-High FOXA1-Absent TNBC
    (Zenodo, 2026-03-06)
}

\author[1]{Eric Robert Lawson}
\affil[1]{
    OrganismCore\\
    \href{mailto:OrganismCore@proton.me}
         {OrganismCore@proton.me}\\
    ORCID:
    \href{https://orcid.org/0009-0002-0414-6544}
         {0009-0002-0414-6544}
}

\date{March 6, 2026}

% ══════════════════════════════════════════════════════════════
\begin{document}
% ══════════════════════════════════════════════════════════════

\maketitle
\thispagestyle{fancy}

% ── PRIORITY NOTICE ───────────────────────────────────────────
\begin{center}
\colorbox{urgentbox}{
\begin{minipage}{0.88\textwidth}
\vspace{0.5em}
\centering
\textbf{PRIORITY ESTABLISHMENT DOCUMENT}\\[0.3em]
\small
As of 2026-03-06, no registered clinical trial exists
anywhere in the world for tazemetostat as maintenance
therapy after neoadjuvant chemotherapy in TNBC
(ClinicalTrials.gov searched 2026-03-05). No solid
tumour adjuvant or maintenance trial for tazemetostat
exists in any indication. This document establishes
priority on the specific maintenance setting, patient
selection criteria, treatment protocol, and primary
endpoint described herein. This proposal is distinct
from CS-LIT-16 (tazemetostat $\rightarrow$ fulvestrant
conversion therapy in active/metastatic TNBC). Both
proposals use the same drug in the same tumour type
but in different settings, with different mechanisms
of action, different patient populations, different
endpoints, and different biological rationales.
\vspace{0.5em}
\end{minipage}
}
\end{center}

\vspace{0.6em}

% ── REPOSITORY POINTER ────────────────────────────────────────
\begin{center}
\colorbox{repobox}{
\begin{minipage}{0.88\textwidth}
\vspace{0.4em}
\centering
\textbf{Full computational record and companion
reasoning artifacts:}\\[0.3em]
\href{https://github.com/Eric-Robert-Lawson/attractor-oncology}
{\texttt{https://github.com/Eric-Robert-Lawson/attractor-oncology}}
\\[0.2em]
\small
Repository ID: 1172048282 \textbar{}
MIT License \textbar{}
Python \textbar{}
Public\\[0.2em]
Source document:
\texttt{Cancer\_Research/BRCA/DEEP\_DIVE/}
\texttt{BRCA\_Cross\_Subtype\_Literature\_Check.md}
(BRCA-S8h, item CS-LIT-17)\\[0.2em]
EZH2 paradox data:
\texttt{.../FOXA1\_EZH2\_RATIO/ratio4.py}
and
\texttt{ratio\_s4\_results/results/ratio\_s4\_log.txt}\\[0.2em]
Companion deposit (CS-LIT-16):
\textit{Tazemetostat Followed by Fulvestrant in
EZH2-High FOXA1-Absent TNBC}
(Zenodo, 2026-03-06)
\vspace{0.4em}
\end{minipage}
}
\end{center}

\vspace{0.8em}

% ── ABSTRACT ──────────────────────────────────────────────────
\begin{abstract}
\noindent
Triple-negative breast cancer (TNBC) presents a
paradox in the clinic: EZH2-high tumours respond
well to neoadjuvant taxane-anthracycline chemotherapy
(higher pathological complete response rates) but
carry substantially worse long-term outcomes
(higher distant relapse-free survival hazard ratio
at 3--5 years). Both arms of this paradox are
confirmed in GSE25066 ($n=508$; pCR $p<0.0001$;
DRFS HR$=1.363$, $p=0.0047$) and the long-window
arm is confirmed in published independent literature
(Frontiers in Oncology 2020, $\sim$4,000 patients;
Fineberg, Montefiore/Albert Einstein: EZH2 IHC
predicts metastatic disease post-neoadjuvant
chemotherapy in TNBC specifically).

The mechanistic explanation is straightforward:
EZH2 drives proliferation, making tumours
chemotherapy-sensitive in the short term. But the
same EZH2 overexpression silences tumour suppression
programmes (CDKN1A, CDKN2A) and maintains the
epigenetic state that enables late metastatic relapse.
EZH2-high TNBC cells that survive chemotherapy ---
whether as minimal residual disease, dormant
micrometastases, or circulating tumour cells ---
retain the EZH2-driven epigenetic programme and
constitute the reservoir for late relapse.

The post-chemotherapy window between treatment
completion and the 3--5 year late-relapse peak is an
intervention window that is currently completely
unused. We propose that tazemetostat (EZH2 inhibitor,
FDA-approved 2020) administered as maintenance therapy
after taxane-anthracycline completion in EZH2-high
TNBC patients targets this residual EZH2-high
population directly, before it drives late metastatic
disease. This proposal was derived from Waddington
landscape attractor geometry before literature review
and is supported by two independent lines of published
evidence confirming both arms of the EZH2 paradox.

The proposed trial: maintenance tazemetostat
800~mg BID beginning 4--6 weeks after neoadjuvant
chemotherapy completion in patients with EZH2-high
TNBC (H-score $\geq150$), regardless of
pathological complete response status. Primary
endpoint: distant relapse-free survival (DRFS)
at 3 years. No registered trial exists for this
indication. Tazemetostat is FDA-approved.
The trial can be designed and submitted immediately.
\end{abstract}

\vspace{0.5em}
\noindent\textbf{Keywords:}
tazemetostat; maintenance therapy; TNBC;
triple-negative breast cancer; EZH2; EZH2 paradox;
late relapse; neoadjuvant chemotherapy;
post-chemotherapy window; minimal residual disease;
distant relapse-free survival; DRFS; adjuvant;
attractor geometry; Waddington landscape; OrganismCore;
GSE25066; breast cancer; precision oncology;
epigenetic maintenance; H3K27me3; PRC2

\newpage
\tableofcontents
\newpage

% ═══════════════════════════════════════════════��══════════════
\section{Document Metadata}
% ══════════════════════════════════════════════════════════════

\begin{center}
\begin{tabular}{ll}
\toprule
\textbf{Field} & \textbf{Value} \\
\midrule
Document ID       & CS-LIT-17-MAINTENANCE-PROPOSAL \\
Series            & BRCA Deep Dive --- Cross-Subtype \\
Type              & Clinical Trial Proposal ---
                    Priority Establishment Document \\
Parent document   & BRCA-S8h (Cross-Subtype Literature
                    Check, item CS-LIT-17) \\
Companion         & CS-LIT-16 (Tazemetostat $\rightarrow$
                    Fulvestrant, Zenodo 2026-03-06) \\
Master artifact   & FOXA1/EZH2 Ratio Master Reasoning
                    Artifact (Zenodo 2026-03-06) \\
Date              & 2026-03-06 \\
Author            & Eric Robert Lawson \\
Organisation      & OrganismCore \\
ORCID             & 0009-0002-0414-6544 \\
Status            & Permanent \\
Repository        &
  \href{https://github.com/Eric-Robert-Lawson/attractor-oncology}
       {Eric-Robert-Lawson/attractor-oncology} \\
Repository ID     & 1172048282 \\
License           & CC BY 4.0 \\
Biological contradictions & 0 \\
Registered maintenance trials & 0 (2026-03-05) \\
\bottomrule
\end{tabular}
\end{center}

% ══════════════════════════════════════════════════════════════
\section{Distinction from CS-LIT-16}
% ══════════════════════════════════════════════════════════════

These two proposals both involve tazemetostat in TNBC
and are derived from the same OrganismCore framework.
They are distinct clinical strategies targeting
different clinical problems with different mechanisms:

\begin{center}
\begin{tabular}{p{3.5cm}p{5cm}p{5cm}}
\toprule
\textbf{Parameter}
& \textbf{CS-LIT-16}
& \textbf{CS-LIT-17 (this document)} \\
\midrule
Setting
& Active/metastatic TNBC;
  locally advanced TNBC
& Post-neoadjuvant chemotherapy;
  adjuvant/maintenance window \\[0.4em]
Disease status
& Measurable disease present
& Minimal residual disease or
  pCR with residual risk \\[0.4em]
Mechanism
& EZH2 inhibition restores
  FOXA1/ESR1, enabling
  fulvestrant to engage the
  restored ER
& EZH2 inhibition eliminates
  the residual EZH2-high
  population that drives
  late metastatic relapse \\[0.4em]
Second drug
& Fulvestrant (required)
& None (tazemetostat alone
  as maintenance) \\[0.4em]
Primary endpoint
& FOXA1 H-score
  re-expression $\geq50$
  at Week~4 biopsy
& Distant relapse-free
  survival (DRFS) at 3 years \\[0.4em]
Target population
& EZH2-high AND
  FOXA1-absent TNBC
& EZH2-high TNBC regardless
  of FOXA1 status \\[0.4em]
Window
& Any time (active disease)
& 4--6 weeks post-chemo
  completion, before the
  3--5 year relapse peak \\[0.4em]
Biological rationale
& Epigenetic conversion:
  silence $\rightarrow$ luminal
  $\rightarrow$ endocrine target
& Residual cell elimination:
  remove the reservoir before
  it reactivates \\
\bottomrule
\end{tabular}
\end{center}

\noindent These strategies are complementary and can
eventually be combined or run in sequence. A patient
who achieves pCR after neoadjuvant chemotherapy might
receive CS-LIT-17 (maintenance tazemetostat to clear
residual EZH2-high cells) and, if they relapse, receive
CS-LIT-16 (tazemetostat $\rightarrow$ fulvestrant to
convert the relapsed tumour). They address different
moments in the disease trajectory.

% ══════════════════════════════════════════════════════════════
\section{The EZH2 Paradox: Foundation of This Proposal}
% ══════════════════════════════════════════════════════════════

\subsection{What the paradox is}

\begin{center}
\colorbox{paradoxbox}{
\begin{minipage}{0.85\textwidth}
\vspace{0.5em}
\textbf{THE EZH2 PARADOX IN TNBC}

\medskip
EZH2-high TNBC tumours show TWO contradictory
prognostic signals depending on the time window:

\medskip
\textbf{Short window (0--18 months):}
EZH2-high $\rightarrow$ HIGH proliferation
$\rightarrow$ chemo-sensitive $\rightarrow$
higher pCR rate $\rightarrow$ appears favourable

\medskip
\textbf{Long window (18 months -- 5+ years):}
EZH2-high $\rightarrow$ epigenetic silencing
of tumour suppressors $\rightarrow$ residual cells
escape $\rightarrow$ late metastatic relapse
$\rightarrow$ WORSE distant relapse-free survival

\medskip
Both arms are real. Both are confirmed.
They are mechanistically unified by the dual
function of EZH2: it drives proliferation (making
cells chemo-sensitive) AND silences CDKN1A/CDKN2A
and other tumour suppressor programmes (making
survivor cells resistant to long-term control).

\medskip
\textbf{The clinical consequence of the paradox:}
Achieving pCR in EZH2-high TNBC is not the same
as achieving durable remission. The cells that
survive chemotherapy in EZH2-high tumours ---
even as minimal residual disease after apparent
pCR on imaging --- retain the EZH2-driven
epigenetic programme. They are the late-relapse
reservoir.

\textbf{The intervention point the paradox reveals:}
The window between chemotherapy completion and
the late-relapse peak (approximately 3--5 years)
is the only time when the residual EZH2-high
population is accessible, minimal in number, and
targetable with a single approved agent.
\vspace{0.5em}
\end{minipage}
}
\end{center}

\subsection{Quantification of both arms in GSE25066}

Both arms of the paradox were quantified in the
OrganismCore framework's Script 4 validation
(ratio4.py) using GSE25066 (Hatzis et al.\
\textit{JAMA} 2011, $n=508$ neoadjuvant chemotherapy
patients). The full results are recorded in:
\texttt{ratio\_s4\_results/results/ratio\_s4\_log.txt}
in the companion repository.

\begin{center}
\begin{tabular}{p{3.5cm}p{4cm}p{5.5cm}}
\toprule
\textbf{Arm} & \textbf{Statistic} & \textbf{Interpretation} \\
\midrule
Short-window\\(pCR rate)
& $p < 0.0001$\\$n = 508$
& Low FOXA1/EZH2 ratio (EZH2-high)
  $\rightarrow$ pCR more likely.
  EZH2-high tumours are
  chemo-sensitive. \\[0.4em]
Long-window\\(DRFS)
& HR $= 1.363$\\$p = 0.0047$\\$n = 508$
& High EZH2 predicts worse
  distant relapse-free survival.
  EZH2-high patients relapse
  later and more frequently. \\[0.4em]
Both arms\\simultaneously
& Confirmed in the same\\dataset
  ($n=508$)
& The same biological variable
  (EZH2 level) predicts opposite
  outcomes at different time
  windows. This is the paradox
  stated as data. \\
\bottomrule
\end{tabular}
\end{center}

\subsection{Independent published confirmation of
the long-window arm}

The long-window arm (EZH2 = poor long-term prognosis)
is not only confirmed in the OrganismCore framework's
computational analysis. It is independently confirmed
in two published sources:

\begin{center}
\begin{tabular}{p{4cm}p{9cm}}
\toprule
\textbf{Source} & \textbf{Finding} \\
\midrule
Frontiers in Oncology 2020
& EZH2/NSD2 axis predicts poor
  prognosis (RFS, OS, DMFS) in
  $\sim$4,000 breast cancer patients.
  EZH2-high is associated with
  worse distant outcomes across
  all subtypes and specifically
  in TNBC. \\[0.4em]
Fineberg et al.\\
(Montefiore/Albert Einstein
Medical Center, HMP Global)
& EZH2 IHC \textit{specifically}
  predicts metastatic disease in
  TNBC \textit{post-neoadjuvant}
  chemotherapy. This is the most
  directly relevant published
  confirmation: the same clinical
  setting (post-neoadjuvant TNBC)
  as the proposed maintenance trial,
  with EZH2 IHC as the predictor. \\
\bottomrule
\end{tabular}
\end{center}

\noindent The Fineberg finding is particularly
significant. It was published by an independent
group, in the specific clinical setting this proposal
targets, using the exact biomarker (EZH2 IHC) this
proposal uses for patient selection, and it confirms
the exact prognostic claim (EZH2 predicts late
metastatic relapse after neoadjuvant chemo in TNBC)
that motivates the intervention window. This is
independent experimental confirmation of the
framework's derivation from geometry.

\subsection{Why the short-window arm does not
contradict this proposal}

A reviewer might ask: if EZH2-high TNBC responds
better to chemotherapy (short-window arm), why add
tazemetostat after chemotherapy? Would tazemetostat
have removed the chemotherapy benefit by suppressing
EZH2-driven proliferation?

The answer is no, for three reasons:

\begin{enumerate}
  \item \textbf{Timing:} The proposed maintenance
        protocol begins 4--6 weeks after
        chemotherapy completion. The chemotherapy
        has already been administered. The pCR (or
        residual disease) status has already been
        determined. The short-window benefit of
        EZH2-high tumours has already been realised
        before tazemetostat is introduced.

  \item \textbf{Different target:} Chemotherapy
        targets actively dividing cells broadly.
        Tazemetostat targets the EZH2 catalytic
        function --- the H3K27me3 deposition that
        silences tumour suppressors. The two
        mechanisms are orthogonal. Tazemetostat
        does not sensitise or desensitise cells
        to taxanes or anthracyclines.

  \item \textbf{The residual population:} After
        chemotherapy, the surviving cells are a
        selected population. They are the cells
        that were not killed by taxane-anthracycline.
        They are disproportionately EZH2-high cells
        with silenced CDKN1A/CDKN2A that allowed
        them to survive despite the cytotoxic
        assault. These are precisely the cells
        tazemetostat targets --- the survivors,
        not the responders.
\end{enumerate}

% ══════════════════════════════════════════════════════════════
\section{The Mechanistic Basis for Maintenance EZH2
Inhibition}
% ══════════════════════════════════════════════════════════════

\subsection{What EZH2 inhibition does to residual
TNBC cells}

Tazemetostat inhibits EZH2 catalytic activity,
preventing H3K27me3 deposition. In residual EZH2-high
TNBC cells after chemotherapy, this has several
mechanistically predicted consequences:

\begin{center}
\colorbox{mechanismbox}{
\begin{minipage}{0.85\textwidth}
\vspace{0.5em}
\textbf{MECHANISM OF TAZEMETOSTAT IN RESIDUAL TNBC CELLS}

\medskip
\textbf{1. De-repression of tumour suppressor programmes:}\\
H3K27me3 marks on CDKN1A (p21) and CDKN2A (p16)
promoters are removed.\\
p21 and p16 re-express.\\
Cell cycle arrest pathways are restored.\\
Residual cells lose their proliferative escape
mechanism.

\medskip
\textbf{2. Induction of differentiation:}\\
Confirmed independently by Schade et al.\
\textit{Nature} 2024:\\
EZH2 inhibition drives TNBC cells toward a more
differentiated, luminal-like state.\\
Differentiated cells are less metastatic,
less invasive, less capable of seeding
distant metastases.

\medskip
\textbf{3. De-repression of FOXA1 and luminal markers:}\\
Confirmed by Toska et al.\ \textit{Nature Medicine}
2017:\\
EZH2 inhibition re-expresses FOXA1 in ER-negative
breast cancer cells.\\
In the maintenance setting, partial FOXA1
re-expression reduces the invasive/metastatic
phenotype of residual cells even if full ER
conversion does not occur.

\medskip
\textbf{4. Reduction of EZH2-dependent survival
signalling:}\\
EZH2 maintains survival in aggressive breast cancer
cells partly through non-catalytic functions
(protein--protein interactions with oncogenic
transcription factors).\\
Tazemetostat disrupts PRC2 complex stability
beyond its catalytic function.
\vspace{0.5em}
\end{minipage}
}
\end{center}

\subsection{The late-relapse mechanism it interrupts}

The 3--5 year late-relapse peak in TNBC is a
well-documented clinical phenomenon. Its molecular
basis in EZH2-high tumours is:

\begin{enumerate}
  \item Chemotherapy achieves pCR or near-pCR.
        Imaging is clear. The patient appears
        disease-free.

  \item Residual EZH2-high cells --- as dormant
        micrometastases, circulating tumour cells,
        or minimal residual disease below imaging
        threshold --- survive in distant sites
        (bone marrow, lung, liver niches).

  \item These cells are held in a quiescent,
        EZH2-maintained epigenetic state. Their
        CDKN1A and CDKN2A are silenced by H3K27me3.
        Their FOXA1/ESR1 are silenced. They are
        not proliferating, so chemotherapy is
        not an option.

  \item At 3--5 years, triggered by microenvironmental
        signals (inflammatory, hormonal, metabolic),
        these cells re-activate. Their silenced tumour
        suppressors remain silenced. They proliferate
        as metastatic disease.

  \item The patient presents with distant relapse
        at 3--5 years, long after the apparent pCR.

  \textbf{Tazemetostat interrupts this sequence
  at step 3:} by removing H3K27me3 marks during
  the dormancy window, it forces re-expression of
  CDKN1A/CDKN2A in dormant residual cells, prevents
  the epigenetic maintenance of the metastatic
  programme, and reduces the residual cell
  population's capacity to re-activate at year 3--5.
\end{enumerate}

\subsection{Why no other approved agent addresses
this window}

\begin{center}
\begin{tabular}{p{3cm}p{4cm}p{5.5cm}}
\toprule
\textbf{Agent} & \textbf{Mechanism}
& \textbf{Why it does not address the window} \\
\midrule
Capecitabine
& Antimetabolite, cytotoxic
& Targets proliferating cells.
  Dormant residual cells are not
  proliferating. No mechanism of
  action against quiescent
  EZH2-high cells. \\[0.3em]
Olaparib
& PARP inhibitor
& Only active in BRCA1/2-mutant
  TNBC ($\sim$15--20\% of cases).
  No mechanism against the EZH2-driven
  epigenetic maintenance programme. \\[0.3em]
Pembrolizumab\\(adjuvant)
& PD-1 checkpoint inhibitor
& Targets immune-mediated killing
  of cancer cells. Dormant
  micrometastases are immune-invisible
  (low PD-L1, low TIL infiltration
  in the dormant state).
  Does not address epigenetic
  maintenance. \\[0.3em]
Tazemetostat
& EZH2 catalytic inhibitor
& Directly targets the molecular
  mechanism maintaining the dormant
  EZH2-high residual cell population.
  Mechanism of action is specific
  to the EZH2-driven epigenetic
  programme. \\
\bottomrule
\end{tabular}
\end{center}

% ══════════════════════════════════════════════════════════════
\section{The Clinical Trial Proposal}
% ══════════════════════════════════════════════════════════════

\subsection{Trial overview}

\begin{center}
\begin{tabular}{ll}
\toprule
\textbf{Parameter} & \textbf{Specification} \\
\midrule
Phase & II, randomised, open-label \\
Design & Tazemetostat vs observation
         post-neoadjuvant chemotherapy \\
Randomisation & 1:1 stratified by pCR status
                (pCR vs non-pCR/RD) and
                EZH2 H-score ($\geq$200 vs
                150--199) \\
Primary endpoint & Distant relapse-free survival
                   (DRFS) at 3 years \\
Key secondary & Event-free survival (EFS),
                OS, iDFS \\
Sample size & $n = 200$--$300$
              (100--150 per arm) \\
Estimated duration & 5 years (2 years accrual,
                     3 years minimum follow-up) \\
Setting & Any oncology centre with IHC
          capability and access to
          tazemetostat \\
Drug approved & Yes (FDA-approved 2020) \\
\bottomrule
\end{tabular}
\end{center}

\subsection{Patient selection criteria}

\subsubsection{Inclusion criteria}

\begin{enumerate}
  \item Histologically confirmed TNBC
        (ER $<1\%$, PR $<1\%$, HER2 IHC 0--1+
        or FISH non-amplified)

  \item EZH2 H-score $\geq 150$ on pre-treatment
        diagnostic biopsy (EZH2 IHC, validated
        antibody; Abcam ab3748 or equivalent)

  \item Completed standard neoadjuvant
        taxane-anthracycline chemotherapy
        (with or without pembrolizumab per
        current standard of care)

  \item Undergone definitive surgery

  \item Any pathological response status
        (pCR or residual disease both eligible)

  \item ECOG performance status 0--2 at
        randomisation

  \item Adequate organ function (hepatic,
        renal, haematological)

  \item Beginning maintenance within 6 weeks
        of surgery/chemotherapy completion

  \item Willing to provide pre-treatment
        biopsy tissue for EZH2 IHC if
        archival tissue is insufficient
\end{enumerate}

\subsubsection{Exclusion criteria}

\begin{enumerate}
  \item EZH2 H-score $< 150$ on pre-treatment
        biopsy (not in the target population)

  \item Evidence of distant metastatic disease
        at randomisation (M1 disease requires
        CS-LIT-16 conversion protocol,
        not maintenance)

  \item Prior EZH2 inhibitor therapy

  \item BRCA1/2 germline mutation with completed
        or planned olaparib adjuvant therapy
        (standard of care conflict)

  \item Active concurrent malignancy

  \item Pregnancy or breast-feeding

  \item Concurrent strong CYP3A4 inhibitors
        or inducers (tazemetostat interaction)

  \item Known EZH2 loss-of-function mutation
        (tazemetostat mechanism requires
        functional EZH2)
\end{enumerate}

\subsection{Stratification and randomisation}

Patients are stratified at randomisation by:

\begin{enumerate}
  \item \textbf{pCR status:} pCR (ypT0/is ypN0)
        vs residual disease. This is the most
        important stratification variable because
        pCR patients have apparent disease
        elimination but retain residual risk
        from the EZH2-high dormant population.
        RD patients have confirmed residual
        disease and unambiguous residual EZH2-high
        cell burden.

  \item \textbf{EZH2 H-score:} $\geq200$ vs
        150--199. Higher EZH2 expression
        predicts greater residual cell burden
        and greater maintenance benefit.
        This stratification allows subgroup
        analysis of dose-response relationship
        between EZH2 level and maintenance
        benefit.

  \item \textbf{Prior pembrolizumab:} yes vs no.
        Pembrolizumab has become standard of
        care in high-risk TNBC neoadjuvant
        regimens. Stratification controls
        for its effect on residual immune
        microenvironment.
\end{enumerate}

\subsection{Treatment protocol}

\begin{center}
\colorbox{lockbox}{
\begin{minipage}{0.85\textwidth}
\vspace{0.5em}
\textbf{TREATMENT PROTOCOL}

\medskip
\textbf{Experimental arm --- Tazemetostat maintenance:}\\
Tazemetostat 800~mg orally twice daily (BID)\\
Continuous dosing\\
Duration: 24 months (2 years) or until:\\
\quad --- Disease recurrence (distant or local)\\
\quad --- Unacceptable toxicity\\
\quad --- Patient withdrawal\\
\quad --- 24-month protocol completion (whichever first)

\medskip
\textit{Dose rationale:} 800~mg BID is the approved
dose for follicular lymphoma and epithelioid sarcoma.
It achieves $>90\%$ EZH2 inhibition at trough
concentrations in published PK/PD data from the
E7438-G000-101 Phase I trial (Italiano et al.\
Cancer Discovery 2018). Full EZH2 inhibition
throughout the maintenance window is required to
prevent H3K27me3 re-accumulation on tumour
suppressor promoters.

\medskip
\textbf{Control arm --- Observation:}\\
Standard post-operative surveillance per local
institutional guidelines.\\
No placebo (open-label design).

\medskip
\textbf{Timing of initiation:}\\
Tazemetostat begins 4--6 weeks after the later of:\\
\quad --- Final chemotherapy administration, or\\
\quad --- Definitive surgery\\
This window allows haematological recovery from
chemotherapy before beginning continuous oral
dosing.

\medskip
\textbf{Dose modifications:}\\
Per published tazemetostat prescribing information
(Tazverik, Epizyme/Ipsen).\\
Grade 3/4 haematological toxicity: dose reduction
to 600~mg BID.\\
Grade 3/4 non-haematological toxicity: hold and
reassess per standard protocol.
\vspace{0.5em}
\end{minipage}
}
\end{center}

\subsection{Endpoints}

\subsubsection{Primary endpoint}

\begin{itemize}
  \item \textbf{Distant relapse-free survival (DRFS)
        at 3 years}

  \item Definition: time from randomisation to
        first distant relapse (bone, lung, liver,
        brain, other distant site) or death from
        any cause, whichever occurs first

  \item 3-year DRFS is chosen rather than 5-year
        because the EZH2 paradox data shows the
        relapse peak in EZH2-high TNBC occurs
        at 3--5 years. The 3-year timepoint
        captures the intervention window directly
        and allows trial completion within a
        feasible regulatory timeline.

  \item The Fineberg data (EZH2 IHC predicts
        metastatic disease post-neoadjuvant in TNBC)
        provides external evidence that this
        endpoint is the correct one for this
        biomarker-selected population.
\end{itemize}

\subsubsection{Key secondary endpoints}

\begin{enumerate}
  \item \textbf{Event-free survival (EFS):} time
        to local or distant recurrence, second
        primary cancer, or death

  \item \textbf{Overall survival (OS):}
        time to death from any cause

  \item \textbf{Invasive disease-free survival
        (iDFS):} per standardised STEEP criteria

  \item \textbf{DRFS at 5 years:}
        captures the full late-relapse window

  \item \textbf{pCR subgroup analysis:}
        DRFS in pCR patients specifically.
        This is the critical subgroup --- patients
        who appeared cured but retain EZH2-high
        residual risk.

  \item \textbf{EZH2 H-score $\geq200$ subgroup:}
        DRFS in the highest-EZH2 patients.
        Expected to show the strongest maintenance
        benefit.
\end{enumerate}

\subsubsection{Exploratory biomarker endpoints}

\begin{enumerate}
  \item \textbf{ctDNA clearance:} serial ctDNA
        at baseline (randomisation), 6 months,
        12 months, 24 months. Hypothesis:
        tazemetostat arm shows faster ctDNA
        clearance due to residual cell
        elimination.

  \item \textbf{EZH2 H-score in residual disease
        (RD patients):} from surgical specimen.
        Correlation with DRFS in the RD subgroup.

  \item \textbf{H3K27me3 in residual disease:}
        pharmacodynamic confirmation (requires
        on-treatment biopsy if patient consents;
        not mandatory).

  \item \textbf{CDKN1A (p21) re-expression:}
        on-treatment biopsy at 3 months (optional,
        consenting patients only). Predicted to
        increase in tazemetostat arm as EZH2
        inhibition de-represses CDKN1A.

  \item \textbf{FOXA1 partial re-expression:}
        on-treatment biopsy at 3 months.
        In the maintenance setting, complete FOXA1
        conversion to luminal identity is not
        required or expected. Partial re-expression
        (H-score increase from $<20$ to 20--100)
        would confirm target engagement in the
        residual population without triggering
        the CS-LIT-16 endocrine conversion
        mechanism.
\end{enumerate}

\subsection{Sample size justification}

The historical 3-year DRFS in EZH2-high TNBC is
estimated at approximately 60--65\% based on:

\begin{itemize}[noitemsep]
  \item DRFS HR $= 1.363$ in GSE25066 for EZH2-high
        vs EZH2-low TNBC patients ($p = 0.0047$)
  \item Published 3-year DRFS for all-comers TNBC
        post-neoadjuvant: approximately 65--70\%
  \item The EZH2-high subgroup is expected to be
        at the lower end ($\sim$60\%) based on
        the HR from GSE25066
\end{itemize}

A clinically meaningful improvement in 3-year
DRFS from 60\% to 72\% (hazard ratio $\approx 0.65$,
consistent with published adjuvant trials in TNBC)
requires:

\begin{center}
\begin{tabular}{lc}
\toprule
\textbf{Parameter} & \textbf{Value} \\
\midrule
Control arm 3-year DRFS & 60\% \\
Experimental arm 3-year DRFS & 72\% \\
Hazard ratio & 0.65 \\
Power & 80\% \\
Two-sided $\alpha$ & 0.05 \\
Required events & $\sim$80 \\
Required patients (accounting for\\
attrition and censoring) & 200--250 \\
Accrual period & 2 years \\
Minimum follow-up & 3 years \\
Total trial duration & 5 years \\
\bottomrule
\end{tabular}
\end{center}

\noindent The EZH2-high (H-score $\geq150$) selection
criterion is expected to enrich for approximately
30--40\% of all TNBC patients based on published
EZH2 IHC frequency data in TNBC cohorts. A centre
treating 50--80 TNBC patients per year could
contribute 15--30 eligible patients per year.
Multi-centre accrual of 5--8 sites achieves the
required sample in the 2-year accrual window.

% ══════════════════════════════════════════════════════════════
\section{The Clinical Gap: Why This Trial Does Not Exist}
% ══════════════════════════════════════════════════════════════

\subsection{ClinicalTrials.gov search results}

Searched 2026-03-05.

\begin{center}
\begin{tabular}{ll}
\toprule
\textbf{Search} & \textbf{Result} \\
\midrule
Tazemetostat adjuvant breast cancer
& No trials \\
Tazemetostat maintenance breast cancer
& No trials \\
Tazemetostat TNBC
& No trials \\
EZH2 inhibitor adjuvant TNBC
& No trials \\
EZH2 inhibitor maintenance solid tumour
& No trials \\
Tazemetostat solid tumour adjuvant
& No trials \\
\midrule
\textbf{Tazemetostat approved indications}
& EZH2-mutant follicular lymphoma (2020)\\
& Epithelioid sarcoma (2020) \\
\textbf{Active tazemetostat trials}
& NCT05023655: ARID1A-mutant\\
& solid tumours (Phase II)\\
& NCI Pediatric MATCH basket \\
& \textbf{None in breast cancer} \\
& \textbf{None in maintenance/adjuvant setting} \\
\bottomrule
\end{tabular}
\end{center}

\subsection{Why the gap exists}

The maintenance setting requires three insights
assembled together that do not appear together
in any published paper:

\begin{enumerate}
  \item The EZH2 paradox exists and has two
        quantifiable arms (confirmed in GSE25066
        and published literature but not named
        or assembled as a single phenomenon
        before this framework).

  \item The post-chemotherapy window between
        apparent pCR and late relapse is the
        specific intervention point the paradox
        reveals (this is the original synthesis
        contribution --- deriving the intervention
        timing from the paradox's two-window
        structure).

  \item Tazemetostat is the correct agent because
        it directly targets the molecular mechanism
        (EZH2-driven H3K27me3) maintaining the
        dormant residual population (this requires
        connecting the EZH2 paradox data to the
        approved EZH2 inhibitor --- a connection
        not made in any published paper).
\end{enumerate}

The three published building blocks exist. The
assembly into a maintenance trial design for TNBC
with a specific patient selection criterion
(EZH2 H-score $\geq150$) and a specific endpoint
(DRFS at 3 years targeting the paradox window)
is the original contribution of this document.

% ══════════════════════════════════════════════════════════════
\section{Relationship to the EZH2 Paradox Framework}
% ═════════════════���════════════════════════════════════════════

\subsection{CS-LIT-10 and CS-LIT-17: same paradox,
different clinical applications}

The EZH2 paradox (CS-LIT-10 in BRCA-S8h) is a
named finding in the OrganismCore framework. It
identifies both arms, quantifies them simultaneously
in one dataset, and derives two clinical consequences:

\begin{center}
\begin{tabular}{p{3cm}p{5cm}p{5cm}}
\toprule
\textbf{Finding} & \textbf{CS-LIT-10}
& \textbf{CS-LIT-17 (this document)} \\
\midrule
What it is
& The EZH2 paradox: naming,
  quantification, and
  mechanistic unification
  of both arms
& The clinical proposal that
  follows from the paradox:
  maintenance tazemetostat
  to close the late-relapse
  window \\[0.3em]
Verdict
& CONVERGENT-NOVEL
& NOVEL \\[0.3em]
Type
& Observation
& Action \\[0.3em]
Status
& Recorded in BRCA-S8h
& This document \\
\bottomrule
\end{tabular}
\end{center}

CS-LIT-10 names the paradox and shows it is real.
CS-LIT-17 specifies what to do about it. They are
logically sequential. Neither alone is sufficient
for clinical translation. Together they constitute
a complete case: here is the problem (paradox),
here is the mechanism (EZH2-driven late relapse),
here is the solution (maintenance tazemetostat),
here is the patient population (EZH2 H-score
$\geq150$ post-neoadjuvant TNBC), here is the
endpoint (DRFS at 3 years).

\subsection{Derivation provenance}

The EZH2 paradox and the maintenance proposal were
derived from Waddington landscape attractor geometry
before the literature check (BRCA-S8h) was performed.
The literature check confirmed the long-window arm
independently from published sources (Frontiers
Oncology 2020, Fineberg). The derivation sequence
in the repository is version-controlled with
cryptographic timestamps:

\begin{center}
\begin{tabular}{lll}
\toprule
\textbf{Document} & \textbf{Content} & \textbf{Status} \\
\midrule
BRCA-S8g  & Script 3 results --- EZH2
             paradox first quantified
           & Geometry derivation \\
BRCA-S8h  & Literature check ---
             CS-LIT-17 confirmed NOVEL
           & Literature confirmed \\
CS-LIT-16 & Companion trial proposal
             (tazemetostat $\rightarrow$
             fulvestrant)
           & Deposited Zenodo
             2026-03-06 \\
CS-LIT-17 & This document
           & Deposited Zenodo
             2026-03-06 \\
\bottomrule
\end{tabular}
\end{center}

All documents are publicly accessible at:
\href{https://github.com/Eric-Robert-Lawson/attractor-oncology}
{\texttt{github.com/Eric-Robert-Lawson/attractor-oncology}}

% ══════════════════════════════════════════════════════════════
\section{Integration with CS-LIT-16: A Combined
Strategy}
% ══════════════════════════════════════════════════════════════

CS-LIT-16 and CS-LIT-17 can operate as a combined
sequential strategy for the same patient:

\begin{center}
\colorbox{mechanismbox}{
\begin{minipage}{0.85\textwidth}
\vspace{0.5em}
\textbf{COMBINED STRATEGY FOR EZH2-HIGH TNBC}

\medskip
\textbf{Line 1 (neoadjuvant):}\\
Standard taxane-anthracycline $\pm$ pembrolizumab

\medskip
\textbf{Line 2 (maintenance/adjuvant) ---
CS-LIT-17:}\\
Tazemetostat 800~mg BID $\times$ 24 months\\
Target: residual EZH2-high dormant cells\\
Goal: prevent late relapse

\medskip
\textbf{If distant relapse occurs (Line 3) ---
CS-LIT-16:}\\
Tazemetostat 800~mg BID $\times$ 4 weeks\\
$\rightarrow$ FOXA1 H-score at week 4 biopsy\\
$\rightarrow$ If FOXA1 $\geq50$: add fulvestrant\\
Target: convert relapsed EZH2-high TNBC to
endocrine-responsive state\\
Goal: engage fulvestrant on the restored ER

\medskip
\textbf{Logic of the combined strategy:}\\
CS-LIT-17 reduces the probability that relapse
occurs (maintenance, prevention).\\
CS-LIT-16 provides a mechanism-based treatment
if relapse occurs despite maintenance (conversion,
rescue).\\
The two proposals are not competing. They are
sequential lines in the same disease trajectory.
\vspace{0.5em}
\end{minipage}
}
\end{center}

This combined strategy does not need to be tested
in a single trial. Each proposal stands alone
and can be tested independently. The combination
is noted here to establish that the two proposals
are part of a coherent framework, not isolated
predictions.

% ══════════════════════════════════════════════════════════════
\section{Locked Priority Statement}
% ══════════════════════════════════════════════════════════════

\begin{center}
\colorbox{urgentbox}{
\begin{minipage}{0.88\textwidth}
\vspace{0.6em}
\textbf{LOCKED AS OF 2026-03-06}

\medskip
\textbf{The proposal:}
Tazemetostat 800~mg BID $\times$ 24 months as
maintenance therapy beginning 4--6 weeks after
neoadjuvant chemotherapy completion in patients
with EZH2-high (H-score $\geq150$) TNBC,
regardless of pCR status.

\medskip
\textbf{The primary endpoint:}
Distant relapse-free survival (DRFS) at 3 years.

\medskip
\textbf{The mechanistic basis:}
The EZH2 paradox: EZH2-high TNBC is chemo-sensitive
short-term but late-relapsing long-term.
The residual EZH2-high population after chemotherapy
maintains dormancy through H3K27me3-mediated
suppression of CDKN1A, CDKN2A, and luminal markers.
Tazemetostat removes those marks during the
dormancy window, before the 3--5 year relapse peak.

\medskip
\textbf{Supporting evidence:}
GSE25066: DRFS HR$=1.363$, $p=0.0047$ ($n=508$).\\
Frontiers Oncology 2020: EZH2 predicts poor prognosis
in $\sim$4,000 breast cancer patients.\\
Fineberg (Montefiore): EZH2 IHC predicts metastatic
TNBC post-neoadjuvant specifically.

\medskip
\textbf{The gap:}
No registered maintenance or adjuvant trial for
tazemetostat exists in any solid tumour indication
as of 2026-03-06. The trial can be designed and
submitted today.

\medskip
\textbf{Derived from:}
OrganismCore Waddington landscape attractor geometry.
Prediction locked before literature review.
Repository:
\href{https://github.com/Eric-Robert-Lawson/attractor-oncology}
{\texttt{github.com/Eric-Robert-Lawson/attractor-oncology}}

\medskip
\hfill --- Eric Robert Lawson, March 6, 2026
\vspace{0.4em}
\end{minipage}
}
\end{center}

% ══════════════════════════════════════════════════════════════
\section*{Repository and Reproducibility}
\addcontentsline{toc}{section}{Repository and
Reproducibility}
% ══════════════════════════════════════════════════════════════

\begin{center}
\href{https://github.com/Eric-Robert-Lawson/attractor-oncology}
{\texttt{https://github.com/Eric-Robert-Lawson/attractor-oncology}}\\[0.3em]
\small Repository ID: 1172048282 \textbar{}
MIT License \textbar{}
Python \textbar{}
Public
\end{center}

\begin{center}
\begin{tabular}{ll}
\toprule
\textbf{File} & \textbf{Path} \\
\midrule
Literature check (parent)
& \texttt{Cancer\_Research/BRCA/DEEP\_DIVE/}\\
& \texttt{BRCA\_Cross\_Subtype\_Literature\_Check.md} \\
EZH2 paradox script
& \texttt{.../FOXA1\_EZH2\_RATIO/ratio4.py} \\
Script 4 full log
& \texttt{.../ratio\_s4\_results/results/}\\
& \texttt{ratio\_s4\_log.txt} \\
Master FOXA1/EZH2 artifact
& \texttt{.../FOXA1\_EZH2\_RATIO/}\\
& \texttt{Master\_Artifact\_After\_Scripts.md} \\
Companion CS-LIT-16
& Zenodo deposit, 2026-03-06 \\
\bottomrule
\end{tabular}
\end{center}

% ══════════════════════════════════════════════════════════════
\section*{Author Statement}
\addcontentsline{toc}{section}{Author Statement}
% ══════════════════════════════════════════════════════════════

This work was conducted independently by Eric Robert
Lawson under the OrganismCore research programme.
No institutional funding was received. No conflicts
of interest exist. No pharmaceutical company was
involved in this proposal. Tazemetostat is manufactured
by Epizyme (now Ipsen). The manufacturer had no
involvement in this work.

This document is a scientific priority establishment
record and trial proposal rationale. It is not a
registered clinical trial protocol. Actual trial
conduct requires ethics committee approval, regulatory
submission, institutional sponsorship, and full
protocol development by qualified clinical
investigators.

% ══════════════════════════════════════════════════════════════
\section*{Contact}
\addcontentsline{toc}{section}{Contact}
% ══════════════════════════════════════════════════════════════

\begin{tabular}{ll}
Author        & Eric Robert Lawson \\
Organisation  & OrganismCore \\
Email         &
  \href{mailto:OrganismCore@proton.me}
       {OrganismCore@proton.me} \\
ORCID         &
  \href{https://orcid.org/0009-0002-0414-6544}
       {https://orcid.org/0009-0002-0414-6544} \\
Repository    &
  \href{https://github.com/Eric-Robert-Lawson/attractor-oncology}
       {github.com/Eric-Robert-Lawson/attractor-oncology} \\
Homepage      &
  \href{https://github.com/Eric-Robert-Lawson/OrganismCore}
       {github.com/Eric-Robert-Lawson/OrganismCore} \\
Date          & 2026-03-06 \\
\end{tabular}

\vspace{2em}

\begin{center}
\itshape
``EZH2-high TNBC responds to chemotherapy.\\
Then it comes back at year three.\\
The window between those two events\\
is empty. Tazemetostat belongs there.''

\medskip
\upshape
--- Eric Robert Lawson, March 6, 2026
\end{center}

% ══════════════════════════════════════════════════════════════
\end{document}
